\chapter{Finance, Saving, and Investment (Text Ch. 23)}

\section{Financial institutions and financial markets}

To study the economics of financial institutions and markets we distinguish between
finance and money, physical and financial capital, capital and investment, and wealth and
saving.

\subsection{Finance and money}

The study of \emph{finance} looks at how households and firms obtain and use financial
resources and how they cope with the risks that arise in this activity. The study of
\emph{money} looks at how households and firms use it, how much of it they hold, how banks
create and manage it, and how its quantity influences the economy (next chapter).

\subsection{Physical capital and financial capital}

\begin{definition}[Physical and financial capital]
\emph{Physical capital} is the tools, instruments, machines, buildings and other items
that have been produced in the past and are used today to produce goods and services.
The funds that firms use to buy physical capital are called \emph{financial capital}.
\end{definition}

\subsection{Capital and investment}

\begin{definition}[Gross and net investment]
\emph{Gross investment} is the total amount spent on purchases of new capital and on
replacing depreciated capital. \emph{Depreciation} is the decrease in the quantity of
capital that results from wear and tear and obsolescence. \emph{Net investment} is the
change in the quantity of capital:
\[
  \text{net investment} = \text{gross investment} - \text{depreciation}.
\]
\end{definition}

\begin{example}
A firm starts the year with \$30 million of capital. During the year its capital
depreciates by \$5 million and it buys \$8 million of new capital. Gross investment is
\$8 million, net investment is $8 - 5 = \$3$ million, and it ends the year with
$30 + 3 = \$33$ million of capital.
\end{example}

\subsection{Wealth and saving}

\begin{definition}[Wealth and saving]
\emph{Wealth} is the value of all the things that people own (a stock). \emph{Saving} is
the amount of income that is not paid in taxes or spent on consumption goods and
services (a flow). Wealth changes through saving and through capital gains or losses:
\[
  \Delta \text{wealth} = \underbrace{(\text{income} - \text{net taxes} - \text{consumption})}_{\text{saving}}
  + \text{capital gains} - \text{capital losses}.
\]
\end{definition}

\subsection{Financial capital markets}

Saving is the source of the funds used to finance investment. These funds are supplied
and demanded in three types of financial markets.
\begin{enumerate}
  \item \textbf{Loan markets}: businesses borrow from banks for working capital;
        households take mortgages and car loans.
  \item \textbf{Bond markets}: a \emph{bond} is a promise to make specified payments on
        specified dates. Governments and large firms raise funds by selling bonds. A
        \emph{mortgage-backed security} entitles its holder to income from a package of
        mortgages.
  \item \textbf{Stock markets}: a \emph{stock} is a certificate of ownership and claim to a
        firm's profits. Stocks are traded on exchanges such as the Toronto Stock Exchange.
\end{enumerate}

\subsection{Financial institutions}

\begin{definition}[Financial institution]
A \emph{financial institution} is a firm that operates on both sides of the markets for
financial capital: it borrows in one market and lends in another.
\end{definition}

Key financial institutions in Canada are
\begin{enumerate}
  \item chartered banks;
  \item trust and loan companies;
  \item credit unions and caisses populaires;
  \item pension funds; and
  \item insurance companies.
\end{enumerate}

\subsection{Insolvency and illiquidity}

\begin{definition}[Net worth, insolvency, illiquidity]
A financial institution's \emph{net worth} is the total market value of what it has lent
minus the market value of what it has borrowed. If net worth is positive the institution
is \emph{solvent} and can remain in business; if it is negative the institution is
\emph{insolvent} and goes out of business. An institution is \emph{illiquid} if it is
solvent but has made long-term loans with funds borrowed short term, and faces a sudden
demand to repay what it has borrowed.
\end{definition}

The 2008 global financial crisis involved both: falling house prices made many U.S.\
institutions insolvent, and fear of insolvency made short-term lending dry up, so solvent
institutions became illiquid. Canadian banks fared relatively well.

\subsection{Interest rates and asset prices}

The interest rate on a financial asset is the interest received expressed as a percentage
of the asset's price, usually per year.

\begin{important}[Asset prices and interest rates move in opposite directions]
For an asset that pays a fixed dollar amount, the higher its price, the lower its
interest rate (yield), and vice versa.
\end{important}

\begin{example}
A bond promises to pay \$1{,}050 one year from now. If its price today is \$1{,}000, its
interest rate is $50/1000 = 5\%$. If the bond's price falls to \$980, the interest rate
rises to $(1050 - 980)/980 = 70/980 \approx 7.1\%$.
Similarly, a share paying a \$5 dividend has a 10\% rate of return at a price of \$50 but
only 5\% at a price of \$100.
\end{example}

\section{The loanable funds market}

\begin{definition}[Market for loanable funds]
The \emph{market for loanable funds} is the aggregate of all the individual financial
markets (loans, bonds, stocks).
\end{definition}

\subsection{Funds that finance investment}

From the circular flow, $Y = C + I + G + X - M$, and households use their income for
consumption, saving and net taxes: $Y = C + S + T$. Equating and rearranging,
\[
  \boxed{I = S + (T - G) + (M - X).}
\]
So investment is financed by three sources:
\begin{enumerate}
  \item \textbf{household saving} $S$;
  \item \textbf{government budget surplus} $T - G$ (negative if the government runs a
        deficit); and
  \item \textbf{borrowing from the rest of the world} $M - X$ (when Canada imports more
        than it exports, it borrows the difference from foreigners).
\end{enumerate}
\emph{National saving} is $S + (T - G)$.

\begin{example}
$S = 300$, $T = 450$, $G = 500$, $X = 400$, $M = 460$ (billions). Then
$I = 300 + (450 - 500) + (460 - 400) = 300 - 50 + 60 = 310$. National saving is
$300 - 50 = 250$, and Canada borrows 60 from abroad.
\end{example}

\subsection{The real interest rate}

\begin{definition}[Nominal and real interest rates]
The \emph{nominal interest rate} is the number of dollars that a borrower pays and a
lender receives in interest in a year, as a percentage of the number of dollars borrowed
and lent. The \emph{real interest rate} is the nominal rate adjusted for the effects of
inflation:
\[
  \text{real interest rate} \approx \text{nominal interest rate} - \text{inflation rate}.
\]
\end{definition}

\begin{example}
A loan has a nominal interest rate of 6\% a year and inflation is 2\%. The real interest
rate is about $6 - 2 = 4\%$. (The exact value is $1.06/1.02 - 1 \approx 3.92\%$.) The real
interest rate is the opportunity cost of loanable funds: the real interest paid on
borrowed funds, or forgone on one's own funds.
\end{example}

\subsection{The demand for loanable funds}

\begin{definition}[Demand for loanable funds]
The \emph{quantity of loanable funds demanded} is the total quantity of funds demanded to
finance investment, the government budget deficit and international investment or
lending during a given period. The \emph{demand for loanable funds} is the relationship
between the quantity demanded and the real interest rate when all other influences on
borrowing plans remain the same.
\end{definition}

Investment depends on the real interest rate and on expected profit. A higher real
interest rate means fewer investment projects are profitable, so the demand curve (DLF)
slopes downward. Other things remaining the same, the greater the expected profit from
new capital, the greater the investment and the greater the demand for loanable funds,
so a rise in expected profit shifts DLF rightward. Expected profit rises in an expansion
and with technological change (new products and processes).

\subsection{The supply of loanable funds}

\begin{definition}[Supply of loanable funds]
The \emph{supply of loanable funds} is the relationship between the quantity of loanable
funds supplied and the real interest rate when all other influences on lending plans
remain the same.
\end{definition}

Saving (the main source of supply) depends on
\begin{enumerate}
  \item the real interest rate (a higher rate makes saving more attractive, so SLF slopes
        upward);
  \item disposable income;
  \item expected future income;
  \item wealth; and
  \item default risk.
\end{enumerate}
An increase in disposable income, a decrease in expected future income, a decrease in
wealth, or a fall in default risk increases saving and shifts the supply of loanable
funds rightward.

\subsection{Equilibrium in the loanable funds market}

The equilibrium real interest rate makes the quantity demanded equal the quantity
supplied. In the short run, shifts in DLF (from swings in expected profit) are larger
than shifts in SLF, so the real interest rate and the quantity of funds are volatile; over
time both curves shift rightward.

\begin{example}[Solving for equilibrium]
Suppose (with $Q$ in billions and $r$ in percent)
\[
  Q_D = 1000 - 100r, \qquad Q_S = 100r - 200.
\]
Setting $Q_D = Q_S$: $1000 - 100r = 100r - 200$, so $200r = 1200$, $r = 6\%$ and
$Q = 1000 - 600 = 400$. If expected profits rise and DLF becomes $Q_D = 1200 - 100r$,
then $1200 - 100r = 100r - 200$ gives $r = 7\%$ and $Q = 500$.
\end{example}

\section{Government in the loanable funds market}

\subsection{A government budget surplus}

A surplus adds to the supply of loanable funds (SLF shifts rightward by the amount of the
surplus). The real interest rate falls, household saving decreases, and the quantity of
funds demanded for private investment increases.

\subsection{A government budget deficit and crowding out}

A deficit adds to the demand for loanable funds (DLF shifts rightward by the amount of
the deficit). The real interest rate rises, household saving increases, and private
investment decreases.

\begin{definition}[Crowding-out effect]
The tendency for a government budget deficit to raise the real interest rate and decrease
private investment is called the \emph{crowding-out effect}.
\end{definition}

\begin{center}
\begin{tikzpicture}[scale=0.75]
  \draw[->] (0,0) -- (9,0) node[right] {Loanable funds (\$ billions)};
  \draw[->] (0,0) -- (0,10) node[above] {Real interest rate (\%)};
  \draw[thick, red] (1,3) -- (7,9) node[right] {SLF};
  \draw[thick, blue] (1,9) -- (7,3) node[right] {DLF (private)};
  \draw[thick, blue, dashed] (2,9) -- (7,4) node[right] {DLF + deficit};
  \fill (4,6) circle (2.5pt);
  \fill (4.5,6.5) circle (2.5pt);
  \draw[dotted] (0,6) node[left] {6} -- (4,6) -- (4,0) node[below] {400};
  \draw[dotted] (0,6.5) node[left] {6.5} -- (4.5,6.5) -- (4.5,0) node[below right] {450};
\end{tikzpicture}
\end{center}

\begin{example}[Crowding out, numerically]
Continue with private demand $Q_D = 1000 - 100r$ and supply $Q_S = 100r - 200$
(equilibrium $r = 6\%$, $Q = 400$, all of it private investment). The government now
runs a deficit of \$100 billion, adding 100 to the demand at every interest rate:
$Q_D = 1100 - 100r$. Then
\[
  1100 - 100r = 100r - 200 \implies r = 6.5\%, \quad Q = 100(6.5) - 200 = 450.
\]
Saving rises from 400 to 450. Private investment is $1000 - 100(6.5) = 350$, down from
400: the deficit of 100 \emph{crowds out} 50 of private investment; the other 50 is
financed by extra saving. Check: $350 + 100 = 450$.
\end{example}

\begin{definition}[Ricardo--Barro effect]
The \emph{Ricardo--Barro effect} (Ricardian equivalence) is the claim that taxpayers,
anticipating that a deficit means higher future taxes, increase their saving by the full
amount of the deficit. Then SLF shifts rightward by the same amount as DLF, the real
interest rate does not change and there is no crowding out. Most economists think the
effect is only partial.
\end{definition}

\section{The global loanable funds market}

Financial capital moves freely across borders, so loanable funds are traded in a global
market and there is a single world real interest rate (adjusted for risk). Canada is a
small part of this market, so it takes the world real interest rate as given.
\begin{itemize}
  \item If, at the world rate, Canadian saving is less than Canadian investment demand,
        Canada borrows the difference from the rest of the world (net foreign borrowing,
        $M - X > 0$).
  \item If saving exceeds investment, Canada lends the difference abroad.
\end{itemize}

\begin{example}
With the domestic curves above and a world real interest rate of 5\%: $Q_S = 300$ and
$Q_D = 500$. Canada invests 500, saves 300 and borrows $500 - 300 = 200$ from the rest
of the world. The domestic interest rate stays at 5\% instead of rising to the 6\%
closed-economy equilibrium.
\end{example}

\section{Summary}

\begin{itemize}
  \item Physical capital is produced inputs; financial capital is the funds used to buy
        it. Net investment $=$ gross investment $-$ depreciation.
  \item Saving is a flow; wealth is a stock that changes by saving plus capital gains.
  \item Funds flow through loan, bond and stock markets and through financial
        institutions; insolvency (negative net worth) differs from illiquidity.
  \item Asset prices and interest rates move inversely.
  \item $I = S + (T - G) + (M - X)$. Real interest rate $\approx$ nominal $-$ inflation.
  \item DLF slopes down (investment), shifts with expected profit; SLF slopes up
        (saving), shifts with disposable income, expected future income, wealth and
        default risk.
  \item A deficit raises the real interest rate and crowds out investment (unless the
        Ricardo--Barro effect is complete). In the global market Canada borrows or lends
        at the world real interest rate.
\end{itemize}
